\subsection{Introduction to the FDG model}
\label{subsec:fdg-intro}
% REVISED (2026-10-24): Theorems Consensus, Disagreement and Corollary Average consensus
% are now stated exactly as in the course slides (only the notation is adapted).

\subsubsection{The FDG model}
\label{subsubsec:fdg-model}

Consider a group of agents $V = \{1, 2, \dots, n\}$, where $x_i \in \mathbb{R}$ is the
opinion of agent $i$ on a given topic, and a weighted digraph
$G = (V, E, \{a_{ij}\}_{(i,j) \in E})$ describing who influences whom, with the
convention of Section~\ref{subsec:prelim}. In the
French--Harary--DeGroot (FDG) model~\cite{baggio_dt}, called French--DeGroot model
in~\cite{proskurnikov2017tutorial}, the new opinion of each agent is a weighted average
of the opinions of the agents it is influenced by,
% REVISED (2026-10-28, review): shortened; the edge convention is already stated in Section 1.1. Was:
% Consider a group of individuals $V = \{1, 2, \dots, n\}$, where $x_i \in \mathbb{R}$
% is the opinion that individual $i$ holds on a given topic. The influences among the
% individuals are described by a weighted digraph
% $G = (V, E, \{a_{ij}\}_{(i,j) \in E})$: following the convention of
% Section~\ref{subsec:prelim}, an edge $(i,j) \in E$, with weight $a_{ij} > 0$, means
% that individual $i$ is influenced by individual $j$.
%
% The French--Harary--DeGroot (FDG) model~\cite{baggio_dt}, called French--DeGroot
% model in~\cite{proskurnikov2017tutorial}, describes how each individual, or agent,
% updates its opinion: the new opinion is a weighted average of the opinions of the
% agents it is influenced by,
\begin{equation}\label{eq:fdg}
  x_i(k+1) = \sum_{j=1}^{n} a_{ij}\, x_j(k), \qquad \sum_{j=1}^{n} a_{ij} = 1,
\end{equation}
or, in vector form, $x(k+1) = A x(k)$, where $A$ is the weighted adjacency matrix
of $G$ (Definition~\ref{def:adjacency}). The constraint on the weights means that
$A$ is row-stochastic (Definition~\ref{def:row-stochastic}). In the terminology
of~\cite{baggio_dt}, the FDG model is a \emph{DT averaging system}.

Since $x(k) = A^k x(0)$, the asymptotic behaviour of the model is determined by the
powers of $A$. Three outcomes are possible:
\begin{itemize}
  \item \emph{no convergence}: $\lim_{k \to +\infty} x(k)$ does not exist for some
        $x(0)$; this happens when $A$ is not semi-convergent
        (Definition~\ref{def:semiconv}), e.g.\ when $A$ has an
        eigenvalue $\mu \neq 1$ with $|\mu| = 1$;
  \item \emph{consensus}: for every $x(0)$, $\lim_{k \to +\infty} x(k) = \alpha \mathbf{1}_n$
        for some $\alpha \in \mathbb{R}$, which depends on
        $x(0)$~\cite{proskurnikov2017tutorial};
% REVISED (2026-10-24): "for every x(0)" added to the consensus outcome, as in
% Definition 14 of the tutorial (consensus is a property of the model); the other two
% outcomes are unchanged (as in the course slides).
  \item \emph{disagreement}: $\lim_{k \to +\infty} x(k) = v$ with
        $v \neq \alpha \mathbf{1}_n$ for every $\alpha \in \mathbb{R}$.
\end{itemize}


\subsubsection{Consensus and social power}
\label{subsubsec:consensus}

The following theorem characterizes when the FDG model reaches a consensus. Its
condition (A3) involves the subgraph of globally reachable nodes, which is always
strongly connected~\cite{baggio_dt}, so that its periodicity is well defined
(Definition~\ref{def:periodic}).
% REVISED (2026-10-28, review): shortened (the justification of the strong connectivity is dropped). Was:
% The following theorem characterizes when the FDG model reaches a consensus. Its
% condition (A3) involves the subgraph of globally reachable nodes. This set of nodes is
% always strongly connected~\cite{baggio_dt}, since a walk between two globally reachable
% nodes only visits nodes that are themselves globally reachable. Hence its periodicity is well defined
% (Definition~\ref{def:periodic}).

\begin{theorem}[Consensus in DT averaging systems~\cite{baggio_dt}]
\label{thm:consensus}
Let $G$ be a weighted digraph with row-stochastic adjacency matrix $A$. The following
statements are equivalent:
\begin{itemize}
  \item[(A1)] the eigenvalue $1$ is simple and all other eigenvalues $\mu$ satisfy
              $|\mu| < 1$;
  \item[(A2)] $A$ is semi-convergent and
              $\lim_{k \to \infty} A^k = \mathbf{1}_n w^\top$ where
              $w \in \mathbb{R}^n$ satisfies $w \geq 0$,
              $\mathbf{1}_n^\top w = 1$, and $w^\top A = w^\top$;
  \item[(A3)] $G$ contains a globally reachable node and the subgraph of globally
              reachable nodes is aperiodic.
\end{itemize}
If any, and therefore all, of the previous conditions are satisfied, then
\begin{itemize}
  \item[(i)]  $w$ is the left dominant eigenvector of $A$ and $w_i > 0$ iff node $i$
              is globally reachable;
  \item[(ii)] the solution of the DT averaging system satisfies
              $\lim_{k \to \infty} x(k) = \big(w^\top x(0)\big)\mathbf{1}_n$.
\end{itemize}
\end{theorem}

\begin{corollary}[Average consensus~\cite{baggio_dt}]
\label{cor:average}
If $A$ is doubly-stochastic and (A3) holds, then
$w = \frac{1}{n}\mathbf{1}_n$ and
$\lim_{k \to +\infty} x(k) = \big(\frac{1}{n}\sum_{i=1}^{n} x_i(0)\big)\, \mathbf{1}_n$.
\end{corollary}

In Theorem~\ref{thm:consensus}, the left dominant eigenvector is the left eigenvector
associated with the dominant eigenvalue $\rho(A) = 1$
(Theorem~\ref{thm:stoch-spectrum})~\cite{baggio_matrix}.

\paragraph{Social power.}
The entry $w_i$ is called the \emph{social power} of agent $i$: it is the weight of
the initial opinion $x_i(0)$ in the final opinion of the group, i.e.\ in the consensus
value~\cite{proskurnikov2017tutorial}
\begin{equation}\label{eq:social-power}
  \lim_{k \to \infty} x_i(k) = w^\top x(0) = \sum_{j=1}^{n} w_j\, x_j(0)
  \qquad \text{for every agent } i.
\end{equation}
This definition requires a consensus: by (A2), all the rows of
$\lim_{k \to \infty} A^k = \mathbf{1}_n w^\top$ are equal to $w^\top$, so a single vector
describes the whole group. Without a consensus, as in Theorem~\ref{thm:disagreement},
the rows of the limit matrix differ, each agent ends at its own weighted average of the
initial opinions, and no single vector plays the role of $w$.
% REVISED (2026-10-28, review): shortened, same content. Was:
% This definition requires a consensus. By (A2), $\lim_{k \to \infty} A^k = \mathbf{1}_n w^\top$: all the rows of the
% limit matrix are equal to $w^\top$, so every agent ends at the same weighted average
% $w^\top x(0)$ of the initial opinions, and a single vector $w$ describes the whole
% group. Without a consensus, as in Theorem~\ref{thm:disagreement}, the limit matrix
% still exists, but its rows are not all equal, otherwise every agent would end at the
% same value for every $x(0)$: each agent ends at its own weighted average of the
% initial opinions, so the weight of $x_j(0)$ depends on the agent, and no single vector
% plays the role of $w$.
% REVISED (2026-10-24): added why the social power requires a consensus (rows of the
% limit matrix all equal to w^T under (A2), not all equal in the disagreement case).

Point (i) has a direct social interpretation: the consensus value is the same for
every agent, so the initial opinion of agent $i$ can enter it only if it reaches
\emph{all} the agents, i.e.\ if $i$ is globally reachable. Hence the consensus value
depends only on the initial opinions of the globally reachable agents.
% REVISED (2026-10-28, review): shortened, same content. Was:
% Point (i) has a direct social interpretation. Since the consensus value
% $w^\top x(0)$ is the same for every agent, the initial opinion of agent $i$ can
% contribute to it only if it reaches \emph{all} the agents. If $i$ is not globally
% reachable, there is at least one agent that is influenced by $i$ neither directly
% nor through intermediaries: the opinion of $i$ never reaches that agent, so it
% cannot enter the common value, and $w_i = 0$. Conversely, the opinion of a globally
% reachable agent reaches the whole group and has a positive weight. Hence the
% consensus value depends only on the initial opinions of the globally reachable
% agents.

Among the globally reachable agents, the social power depends on their position in the
network, and it can be regarded as a centrality measure~\cite{proskurnikov2017tutorial}.
For instance, in the model of French with $n = 3$ agents and
\[
  A = \begin{bmatrix} 1/2 & 1/2 & 0 \\ 1/3 & 1/3 & 1/3 \\ 0 & 1/2 & 1/2 \end{bmatrix},
\]
all agents are globally reachable, but agent~$2$, the only one listened to by all the
others, has the largest social power, $w = (2/7,\, 3/7,\, 2/7)^\top$~\cite{proskurnikov2017tutorial}.
% REVISED (2026-10-28, review): shortened (the description of who listens to whom can be read on the matrix). Was:
% agents~$1$ and~$3$ listen to themselves and to agent~$2$, while agent~$2$ listens to
% everybody. All agents are globally reachable, but agent~$2$, the only one listened to by
% all the others, has the largest social power:
% $w = (2/7,\, 3/7,\, 2/7)^\top$~\cite{proskurnikov2017tutorial}.


\subsubsection{Disagreement}
\label{subsubsec:disagreement}

Theorem~\ref{thm:consensus} requires a single globally reachable group, i.e.\ a
unique sink in the condensation digraph $C(G)$. We now consider the case in which
$C(G)$ has $n_s \geq 2$ sinks. Each sink is a strongly connected component, so its
periodicity is well defined (Definition~\ref{def:periodic}).

\begin{theorem}[Disagreement in DT averaging systems~\cite{baggio_dt}]
\label{thm:disagreement}
Let $G$ be a weighted digraph with row-stochastic adjacency matrix $A$. Assume that
the condensation digraph $C(G)$ has $n_s \geq 2$ sinks. The following statements are
equivalent:
\begin{itemize}
  \item[(A1)] the eigenvalue $1$ is semi-simple with multiplicity $n_s$ and all other
              eigenvalues $\mu$ satisfy $|\mu| < 1$;
  \item[(A2)] $A$ is semi-convergent;
  \item[(A3)] each sink of $C(G)$ is aperiodic.
\end{itemize}
If any, and therefore all, of the previous conditions are satisfied, then the solution
of the DT averaging system satisfies
\[
  \lim_{k \to \infty} x(k) = v \neq \alpha \mathbf{1}_n, \quad \alpha \in \mathbb{R}.
\]
\end{theorem}

The final opinions can be described group by group. The agents of a sink are not
influenced by anyone outside it, so each sink evolves as an FDG model on its own and,
being aperiodic, reaches an internal consensus that depends only on the initial
opinions of its members; every agent outside the sinks converges to a convex
combination of the consensus values of the sinks~\cite{baggio_dt}. In social terms, a
sink is a closed group that listens only to itself, and the agents in between settle
on an intermediate opinion: the disagreement is bounded by the smallest and the largest
consensus value of the sinks.
% REVISED (2026-10-28, review): two paragraphs merged, same content. Was:
% The final opinions can be described group by group. The agents of a sink are not
% influenced by anyone outside the sink, so each sink evolves as an FDG model on its own
% and, being aperiodic, reaches an internal consensus that depends only on the initial
% opinions of its members. Every agent outside the sinks converges instead to a convex
% combination of the consensus values of the sinks~\cite{baggio_dt}.
%
% In social terms, a sink is a closed group that listens only to itself, and two closed
% groups never influence each other: each of them keeps its own consensus value. The
% agents in between listen, directly or through intermediaries, to several closed
% groups and settle on an intermediate opinion between their consensus values. As a
% consequence, the disagreement is bounded: every final opinion lies between
% the smallest and the largest consensus value of the sinks.


\subsubsection{Rosetta stone}
\label{subsubsec:rosetta}

Table~\ref{tab:rosetta}, the so-called \emph{Rosetta stone}~\cite{baggio_dt},
summarizes all the possible behaviours of the FDG model: it translates each property
of the row-stochastic matrix $A$ into a property of the digraph $G$ and into the
asymptotic outcome, with an example for each case; the translation between $A$ and $G$
comes from Theorems~\ref{thm:irreducible} and~\ref{thm:primitive} in cases (i)--(ii),
and from Theorems~\ref{thm:consensus} and~\ref{thm:disagreement} in cases
(iii)--(iv). In case~(i) both agents are
globally reachable and $A$ is doubly-stochastic, so by Corollary~\ref{cor:average} the
group agrees on the average of the initial opinions. In case~(iii) agent~$1$ is
stubborn (Definition~\ref{def:stubborn}) and is the only globally reachable agent:
$w = (1, 0)^\top$ and the group agrees on $x_1(0)$. In case~(iv) both agents are
stubborn: each of them is an aperiodic sink of $C(G)$, and every agent keeps its
initial opinion (Theorem~\ref{thm:disagreement}).
% REVISED (2026-10-28, review): shortened; cases (iii) and (iv) are now read as one and two stubborn agents, the thread of Section 2.2; Theorems 1.12 and 1.13, never cited before, are now cited for cases (i)-(ii). Was:
% The results of this section, together with the consensus and disagreement theorems
% above, classify all the possible behaviours of the FDG model. The classification is
% summarized in Table~\ref{tab:rosetta}, the so-called \emph{Rosetta stone}~\cite{baggio_dt}, which
% translates each property of the row-stochastic matrix $A$ into a property of the
% digraph $G$ and into the corresponding asymptotic outcome; for each case, the table
% also reports an example digraph, its matrix and its spectrum.
%
% The two consensus cases illustrate Theorem~\ref{thm:consensus}. In case~(i) both
% agents are globally reachable and $A$ is doubly-stochastic, so by
% Corollary~\ref{cor:average} $w = (1/2, 1/2)^\top$ and the group agrees on the average
% $\big(x_1(0) + x_2(0)\big)/2$. In case~(iii) agent~$2$ is influenced by agent~$1$
% but not vice versa, so only agent~$1$ is globally reachable: $w = (1, 0)^\top$ and
% the group agrees on $x_1(0)$, whatever the initial opinion of agent~$2$.
%
% Case~(iv) illustrates Theorem~\ref{thm:disagreement} in its simplest form: each agent
% is influenced only by itself, so each agent is a sink of $C(G)$ ($n_s = 2$), both sinks
% are aperiodic thanks to the self-loops, and every agent keeps its initial opinion,
% $x(k) = x(0)$ for all $k$.

\begin{table}[H]
  \centering
  \small
  \renewcommand{\arraystretch}{1.6}
  \begin{tabular}{@{}c c c c >{\raggedright\arraybackslash}m{4.4cm} l@{}}
    \toprule
    Case & Example digraph & $A$ & $\operatorname{spec}(A)$ & Property of $A$ / of $G$ & Outcome \\
    \midrule
    (i)
      & \begin{tikzpicture}[baseline=-0.6ex]
          \node[agent] (a) at (0,0) {1}; \node[agent] (b) at (1.4,0) {2};
          \draw[edge] (a) to[bend left=25] node[weight, above] {1/2} (b);
          \draw[edge] (b) to[bend left=25] node[weight, below] {1/2} (a);
          \draw[selfloop] (a) to node[weight, above] {1/2} (a);
          \draw[selfloop] (b) to node[weight, above] {1/2} (b);
        \end{tikzpicture}
      & $\left[\begin{smallmatrix} 1/2 & 1/2 \\ 1/2 & 1/2 \end{smallmatrix}\right]$
      & $\{1,\, 0\}$
      & $A$ primitive / $G$ strongly connected and aperiodic
      & consensus \\
    (ii)
      & \begin{tikzpicture}[baseline=-0.6ex]
          \node[agent] (a) at (0,0) {1}; \node[agent] (b) at (1.4,0) {2};
          \draw[edge] (a) to[bend left=25] (b);
          \draw[edge] (b) to[bend left=25] (a);
        \end{tikzpicture}
      & $\left[\begin{smallmatrix} 0 & 1 \\ 1 & 0 \end{smallmatrix}\right]$
      & $\{1,\, -1\}$
      & $A$ irreducible, not primitive / $G$ strongly connected and periodic
      & no convergence \\
    (iii)
      & \begin{tikzpicture}[baseline=-0.6ex]
          \node[agent] (a) at (0,0) {1}; \node[agent] (b) at (1.4,0) {2};
          \draw[edge] (b) to node[weight, below] {1/2} (a);
          \draw[selfloop] (a) to (a);
          \draw[selfloop] (b) to node[weight, above] {1/2} (b);
        \end{tikzpicture}
      & $\left[\begin{smallmatrix} 1 & 0 \\ 1/2 & 1/2 \end{smallmatrix}\right]$
      & $\{1,\, 1/2\}$
      & $1$ simple and strictly dominant, $A$ not primitive / $C(G)$ has a unique sink, which is aperiodic
      & consensus \\
    (iv)
      & \begin{tikzpicture}[baseline=-0.6ex]
          \node[agent] (a) at (0,0) {1}; \node[agent] (b) at (1.4,0) {2};
          \draw[selfloop] (a) to (a);
          \draw[selfloop] (b) to (b);
        \end{tikzpicture}
      & $\left[\begin{smallmatrix} 1 & 0 \\ 0 & 1 \end{smallmatrix}\right]$
      & $\{1,\, 1\}$
      & $1$ semi-simple but not simple, strictly dominant / $C(G)$ has at least two sinks, all aperiodic
      & disagreement \\
    (v)
      & \begin{tikzpicture}[baseline=-0.6ex]
          \node[agent] (a) at (0,0) {1}; \node[agent] (b) at (1.1,0) {2};
          \node[agent] (c) at (2.2,0) {3};
          \draw[edge] (a) to[bend left=25] (b);
          \draw[edge] (b) to[bend left=25] (a);
          \draw[selfloop] (c) to (c);
        \end{tikzpicture}
      & $\left[\begin{smallmatrix} 0 & 1 & 0 \\ 1 & 0 & 0 \\ 0 & 0 & 1 \end{smallmatrix}\right]$
      & $\{1,\, 1,\, -1\}$
      & $1$ not strictly dominant / $G$ not strongly connected, at least one sink of $C(G)$ is periodic
      & no convergence \\
    \bottomrule
  \end{tabular}
  \caption{Rosetta stone of the FDG model $x(k+1) = A x(k)$, with $A$ row-stochastic:
  for each case, an example digraph (unlabelled edges have weight $1$), its matrix
  $A$, the spectrum of $A$, the general property of $A$ and of $G$, and the outcome. Cases (i)--(iv)
  and the examples follow~\cite{baggio_dt}; the properties of case (v),
  left blank in the slides, follow from~\cite{proskurnikov2017tutorial}.}
  \label{tab:rosetta}
\end{table}
